\chapter{Mở đầu}

\section{Bối cảnh chuyển đổi số trong lĩnh vực F\&B}

Lĩnh vực F\&B có đặc thù vận hành khác với nhiều loại hình bán lẻ thông thường. Một ca phục vụ trong nhà hàng/quán ăn không kết thúc ngay tại thời điểm khách chọn sản phẩm, mà thường kéo dài qua nhiều bước: khách vào bàn, xem thực đơn, gọi món nhiều lần, nhân viên xác nhận, bếp/bar xử lý từng món, món được phục vụ, hóa đơn được tổng hợp và thanh toán cuối phiên. Trong giờ cao điểm, chỉ cần một mắt xích chậm hoặc sai lệch trạng thái giữa bàn, nhân viên, bếp và thu ngân cũng có thể làm giảm chất lượng trải nghiệm của khách và tăng gánh nặng cho đội ngũ vận hành.

Bối cảnh thị trường F\&B Việt Nam trong những năm gần đây cũng cho thấy nhu cầu tối ưu vận hành, kiểm soát chi phí và tăng hiệu quả phục vụ ngày càng rõ. Các báo cáo và bài viết thị trường về F\&B Việt Nam năm 2025 nhấn mạnh xu hướng phục hồi, cạnh tranh và nhu cầu số hóa quy trình để nâng cao hiệu quả quản lý \cite{ipos-nestle-fnb-report-2025,baochinhphu-fnb-growth-2025}. Đối với một nhà hàng nhỏ, số hóa có thể bắt đầu từ thực đơn điện tử hoặc ghi nhận hóa đơn. Đối với một hệ thống nhiều cửa hàng hoặc một nền tảng phục vụ nhiều nhà hàng, số hóa cần đi xa hơn: dữ liệu thực đơn, bàn, nhân sự, đơn hàng, bếp, thanh toán, báo cáo và gói dịch vụ phải được quản lý thống nhất.

POS (\textit{Point of Sale}) giữ vai trò trung tâm trong quá trình đó. Ở mức đơn giản, POS giúp ghi nhận giao dịch và thanh toán. Trong bối cảnh F\&B, POS còn là nơi nối nhiều trạng thái vận hành: bàn nào đang có khách, đơn nào đang chờ xác nhận, món nào đang được bếp xử lý, hóa đơn nào đang chờ thanh toán, phương thức thanh toán nào đã được ghi nhận và dữ liệu nào cần đưa vào báo cáo cuối ngày. Khi POS được triển khai theo mô hình SaaS, hệ thống không chỉ phục vụ một nhà hàng đơn lẻ mà còn phải quản lý nhiều đơn vị thuê bao, nhiều gói dịch vụ, nhiều cấu hình thanh toán và nhiều phạm vi dữ liệu tách biệt.

Đặt món qua mã QR là một hướng số hóa phù hợp với đặc thù phục vụ tại bàn. Khách có thể quét mã QR, xem thực đơn, thêm món và gửi yêu cầu mà không phụ thuộc hoàn toàn vào việc nhân viên đến bàn ở đúng thời điểm. Tuy nhiên, mã QR chỉ là điểm vào của trải nghiệm. Để trở thành một phần của hệ thống POS, luồng QR ordering cần biết đúng nhà hàng, đúng bàn, đúng phiên phục vụ, đúng trạng thái giỏ, đúng hóa đơn và đúng kênh điều phối bếp. Nếu không được tích hợp với POS, KDS và thanh toán, QR ordering dễ trở thành một thực đơn điện tử tách rời thay vì một năng lực vận hành hoàn chỉnh.

Thanh toán điện tử, đặc biệt là thanh toán qua mã QR/chuyển khoản nhanh, là một yếu tố thực tiễn quan trọng trong bối cảnh Việt Nam. NAPAS mô tả VietQR trong dịch vụ chuyển tiền nhanh 24/7 như hạ tầng mã QR hỗ trợ thông tin chuyển tiền ngân hàng \cite{napas-fastfund-vietqr}. Khi đưa vào hệ thống POS, thanh toán QR không chỉ là hiển thị mã để khách chuyển khoản; hệ thống còn phải gắn giao dịch với hóa đơn, ghi nhận trạng thái thanh toán, xử lý webhook, chống ghi nhận trùng lặp và hỗ trợ đối soát. Điều này làm bài toán thanh toán trở thành một phần của thiết kế hệ thống phân tán, không chỉ là một màn hình giao diện.

Từ các yếu tố trên, đề tài QRTable được đặt trong bối cảnh xây dựng một nền tảng POS theo mô hình SaaS cho F\&B, tích hợp đặt món qua mã QR, điều phối bếp, thanh toán VietQR/SePay, quản trị đơn vị thuê bao và báo cáo vận hành. Kiến trúc microservices được sử dụng như một hướng tổ chức hệ thống theo các miền nghiệp vụ rõ ràng, giúp tách trách nhiệm giữa thực đơn, đơn hàng, bếp, thanh toán, SaaS và truy cập người dùng. Luận văn không xem microservices là lựa chọn mặc nhiên cho mọi hệ thống F\&B; đây là quyết định kiến trúc được phân tích trong phạm vi bài toán QRTable, nơi nhiều miền nghiệp vụ có dữ liệu, vòng đời và quy tắc riêng.

\section{Lý do chọn đề tài}

Quy trình gọi món thủ công trong nhà hàng thường phụ thuộc vào nhân viên như một điểm trung gian giữa khách, POS và bếp. Khi lượng khách tăng, nhân viên phải đồng thời nhận yêu cầu, nhập đơn, xác nhận món, hỏi lại thay đổi, chuyển thông tin cho bếp và xử lý thanh toán. Các thao tác này dễ phát sinh độ trễ, sai sót nhập liệu hoặc thiếu đồng bộ trạng thái. Với mô hình phục vụ tại bàn, một bàn cũng có thể gọi thêm món nhiều lần trong cùng phiên, khiến việc theo dõi đơn, hóa đơn và trạng thái món trở nên phức tạp hơn một giao dịch bán lẻ đơn lẻ.

Các hệ thống POS và nền tảng F\&B tại Việt Nam đã thể hiện nhu cầu thực tế về QR ordering, quản lý nhà hàng, thanh toán và báo cáo. KiotViet có tài liệu hướng dẫn chức năng gọi món qua mã QR \cite{kiotviet-qr-order-doc}; Sapo mô tả phần mềm quản lý nhà hàng với các năng lực quản lý bán hàng, thực đơn, nhân viên và báo cáo \cite{sapo-fnb-restaurant-pos}; iPOS giới thiệu giải pháp O2O/QR ordering để tối ưu chi phí vận hành và tốc độ order \cite{ipos-o2o-qr-order}. Các nguồn này cho thấy bài toán không chỉ có ý nghĩa kỹ thuật mà còn gắn với nhu cầu sản phẩm thực tế. Trong khóa luận này, các hệ thống đó được dùng như bối cảnh và hệ thống liên quan, không dùng để xếp hạng hơn/kém giữa QRTable và sản phẩm thương mại khi không có khảo sát hoặc đo kiểm so sánh trực tiếp.

Điểm đáng nghiên cứu nằm ở việc tích hợp nhiều luồng thành một nền tảng SaaS POS có ranh giới kỹ thuật rõ. Một hệ thống chỉ hiển thị thực đơn QR không giải quyết đầy đủ bài toán vận hành nếu nhân viên vẫn phải tự chuyển thông tin sang POS, bếp không nhận được phiếu kịp thời, thanh toán không gắn với hóa đơn hoặc dữ liệu các nhà hàng dùng chung nền tảng không được cô lập. QRTable vì vậy tập trung vào chuỗi nghiệp vụ từ khách quét QR tại bàn đến giỏ dùng chung, gửi đơn, xác nhận nhân viên, KDS, thanh toán và quản trị tenant.

Mô hình SaaS làm bài toán rộng hơn so với một ứng dụng POS cài đặt riêng. Nền tảng phải hỗ trợ vòng đời đơn vị thuê bao, gói dịch vụ, hạn mức sử dụng, phân quyền theo vai trò, cấu hình thanh toán và khả năng quản trị nền tảng. Chủ nhà hàng, quản lý, nhân viên, bếp/bar, khách hàng và quản trị nền tảng có cách truy cập khác nhau. Các vai trò này không thể dùng chung một cơ chế quyền đơn giản; hệ thống cần phân biệt xác thực nhân viên/chủ/quản trị, phiên khách qua QR, tenant isolation và quyền theo gói dịch vụ.

Kiến trúc microservices được chọn vì bài toán QRTable có nhiều miền nghiệp vụ có trách nhiệm dữ liệu khác nhau. Catalog quản lý thực đơn, bàn, mã QR và tồn kho; Order quản lý phiên, giỏ, đơn và hóa đơn; Kitchen quản lý bản chiếu vận hành cho KDS; Payment quản lý cấu hình và giao dịch thanh toán; SaaS quản lý tenant, gói, thuê bao và hóa đơn gói; User-Access và Authorizer quản lý người dùng, vai trò và định danh. Việc tách các miền này giúp khóa luận phân tích rõ data ownership, service boundary, giao tiếp đồng bộ/bất đồng bộ và các bất biến như idempotency, RBAC, tenant isolation và Saga đại diện.

Vì vậy, đề tài được chọn không nhằm chứng minh sự mới lạ của mã QR hay khẳng định microservices luôn tốt hơn monolith. Đóng góp chính nằm ở việc nghiên cứu, thiết kế và xây dựng một hệ thống phần mềm hoàn chỉnh cho bối cảnh SaaS POS F\&B, trong đó QR ordering được nối với POS, KDS, thanh toán và quản trị tenant bằng kiến trúc có ranh giới rõ.

\section{Phát biểu bài toán}

Bài toán của khóa luận là nghiên cứu và xây dựng nền tảng POS theo mô hình SaaS cho lĩnh vực F\&B, tích hợp đặt món qua mã QR tại bàn và tổ chức hệ thống theo kiến trúc microservices. Hệ thống cần phục vụ đồng thời hai nhóm trải nghiệm: trải nghiệm của khách hàng khi quét QR để gọi món và trải nghiệm của đội ngũ vận hành khi quản lý thực đơn, bàn, đơn hàng, bếp, thanh toán, báo cáo và gói dịch vụ.

Ở phía khách hàng, hệ thống cần cho phép khách quét mã QR tại bàn, truy cập đúng thực đơn của nhà hàng, tham gia phiên phục vụ, thao tác giỏ món dùng chung, gửi đơn chờ nhân viên xác nhận và theo dõi trạng thái đơn. Luồng này phải xử lý được ngữ cảnh tenant/bàn/phiên, tránh tạo đơn trùng lặp và cung cấp cập nhật trạng thái phù hợp mà không yêu cầu khách tạo tài khoản đăng nhập.

Ở phía vận hành nhà hàng, hệ thống cần hỗ trợ chủ nhà hàng, quản lý, nhân viên phục vụ và bếp/bar. Chủ nhà hàng và quản lý cần cấu hình thực đơn, bàn, mã QR, nhân sự, thanh toán, gói thuê bao và báo cáo. Nhân viên cần xem phiên bàn, xác nhận đơn, xử lý hóa đơn và ghi nhận thanh toán. Bếp/bar cần nhận phiếu bếp, xem thứ tự ưu tiên, cập nhật trạng thái món và giúp POS theo dõi món sẵn sàng phục vụ.

Ở phía nền tảng SaaS, hệ thống cần hỗ trợ quản trị viên nền tảng trong việc tạo và quản lý đơn vị thuê bao, gói dịch vụ, hóa đơn gói, trạng thái tenant, quyền truy cập và phân tích nền tảng. Bài toán này yêu cầu cô lập dữ liệu giữa các tenant, kiểm soát quyền theo vai trò và theo gói, đồng thời đảm bảo các dịch vụ không truy cập chéo dữ liệu trái với ranh giới sở hữu.

Về mặt kiến trúc, hệ thống cần tách các miền nghiệp vụ thành các dịch vụ có trách nhiệm rõ, dùng giao tiếp đồng bộ cho các truy vấn/lệnh cần phản hồi trực tiếp và dùng sự kiện bất đồng bộ cho các tác dụng phụ phù hợp. Các bất biến như tenant isolation, idempotency, data ownership, RBAC, entitlement theo gói và consistency giữa các dịch vụ phải được thể hiện trong thiết kế, mã nguồn, kiểm thử và các minh chứng đánh giá.

Bảng~\ref{tab:chapter1-problem-solution} tóm tắt các nhóm vấn đề vận hành chính và hướng giải quyết trong QRTable.

{\footnotesize\sloppy
\setlength{\tabcolsep}{3pt}
\begin{longtable}{>{\raggedright\arraybackslash}p{0.2\textwidth}>{\raggedright\arraybackslash}p{0.34\textwidth}>{\raggedright\arraybackslash}p{0.36\textwidth}}
  \caption{Tóm tắt vấn đề và hướng giải quyết của QRTable.}
  \label{tab:chapter1-problem-solution}\\
  \toprule
  Nhóm vấn đề & Biểu hiện trong vận hành F\&B & Hướng giải quyết trong QRTable \\
  \midrule
  \endfirsthead
  \toprule
  Nhóm vấn đề & Biểu hiện trong vận hành F\&B & Hướng giải quyết trong QRTable \\
  \midrule
  \endhead
  \bottomrule
  \endfoot
  Gọi món tại bàn & Khách phụ thuộc vào thời điểm nhân viên tiếp cận bàn; gọi thêm món trong giờ cao điểm dễ bị chậm hoặc nhập thiếu. & Customer PWA theo phiên QR, thực đơn theo tenant, giỏ dùng chung, phiên bản giỏ và gửi đơn có idempotency. \\
  \addlinespace
  Điều phối bếp & Thông tin từ POS sang bếp/bar dễ đứt đoạn nếu không có trạng thái phiếu bếp rõ ràng. & KDS, phiếu bếp theo trạm, trạng thái dòng món, Kafka event và WebSocket hint/refetch. \\
  \addlinespace
  Thanh toán và đối soát & Hóa đơn cần gắn với đơn, phiên bàn và phương thức thanh toán; webhook/chuyển khoản có thể phát sinh lặp. & Bill/payment lifecycle, cash/VietQR/SePay, webhook idempotency, nhật ký kiểm toán và cầu nối Payment--Order. \\
  \addlinespace
  Quản trị nhiều nhà hàng & Nhiều tenant dùng chung nền tảng có nguy cơ trộn dữ liệu, quyền và hạn mức nếu thiếu ranh giới. & Tenant isolation, RBAC, plan entitlement, lifecycle tenant/subscription và phân tích nền tảng cho Super Admin. \\
\end{longtable}
}

\section{Mục tiêu nghiên cứu và mục tiêu xây dựng hệ thống}

Khóa luận có ba nhóm mục tiêu chính: mục tiêu nghiên cứu, mục tiêu xây dựng hệ thống và mục tiêu đánh giá.

Thứ nhất, về mục tiêu nghiên cứu, khóa luận tổng hợp cơ sở lý thuyết và bối cảnh thực tiễn liên quan đến POS F\&B, đặt món qua mã QR, SaaS và multi-tenancy, kiến trúc microservices, event-driven architecture, consistency/idempotency, giao tiếp thời gian thực và bảo mật trong hệ thống SaaS POS. Các nội dung này không được trình bày như lý thuyết tách rời, mà được dùng để giải thích các quyết định thiết kế của QRTable ở những chương sau. Ví dụ, lý thuyết về multi-tenancy làm nền cho tenant isolation; lý thuyết về microservices làm nền cho service boundary và data ownership; lý thuyết về idempotency và Saga làm nền cho xử lý gửi đơn, webhook thanh toán và phối hợp liên dịch vụ.

Thứ hai, về mục tiêu xây dựng hệ thống, khóa luận hướng đến hiện thực các luồng vận hành cốt lõi của QRTable. Các luồng này gồm: khởi tạo tenant và chủ sở hữu ban đầu; quản lý thực đơn, bàn, khu vực và mã QR; khách tham gia phiên qua QR, thao tác giỏ và gửi đơn; nhân viên xác nhận đơn và bảo toàn tồn kho; KDS nhận phiếu bếp và cập nhật trạng thái món; ghi nhận thanh toán tiền mặt và VietQR/SePay; quản lý nhân sự/quyền; bảng điều khiển và báo cáo theo gói dịch vụ; quản trị nền tảng cho tenant, gói và phân tích tổng quan.

Thứ ba, về mục tiêu đánh giá, khóa luận cần chỉ ra hệ thống đã đáp ứng mục tiêu đến đâu dựa trên bằng chứng có thể truy vết. Việc đánh giá không dừng ở việc hệ thống có client chạy được, mà còn đối chiếu yêu cầu, thiết kế, mã nguồn, kiểm thử, sơ đồ, ảnh minh họa và hiện vật triển khai production. Các kết luận được phân biệt theo loại bằng chứng: đã được kiểm chứng bằng kiểm thử tự động hoặc minh chứng thực thi; được hỗ trợ bởi thiết kế/kiến trúc và mã nguồn; hoặc thuộc hướng phát triển cần thêm bằng chứng. Cách đánh giá này giúp khóa luận giữ được tính học thuật và tránh kết luận vượt quá cơ sở kiến trúc.

Các mục tiêu trên dẫn đến một yêu cầu xuyên suốt: QRTable phải được trình bày như một hệ thống được nghiên cứu, thiết kế, xây dựng và đánh giá có kiểm soát. Luận văn không chỉ mô tả chức năng người dùng, mà còn giải thích vì sao hệ thống được chia thành các dịch vụ, dịch vụ nào sở hữu dữ liệu nào, giao tiếp nào dùng REST/TCP, giao tiếp nào dùng Kafka/WebSocket, và bất biến nào phải được bảo vệ ở từng ranh giới.

\section{Đối tượng và phạm vi nghiên cứu}

Đối tượng nghiên cứu của khóa luận là nền tảng SaaS POS cho nhà hàng/quán F\&B có tích hợp đặt món qua mã QR tại bàn. Trọng tâm không nằm ở mã QR như một kỹ thuật mã hóa dữ liệu đơn lẻ, mà ở quy trình phần mềm phía sau mã QR: định danh tenant và bàn, tạo/duy trì phiên phục vụ, đồng bộ giỏ, gửi đơn, xác nhận đơn, điều phối bếp, thanh toán và quản trị dữ liệu nhiều tenant.

Về phạm vi nghiệp vụ, QRTable tập trung vào các năng lực cần thiết để vận hành một chuỗi quy trình F\&B cơ bản nhưng đầy đủ: quản lý tenant và gói thuê bao; quản lý thực đơn, danh mục, món, bàn, khu vực và mã QR; phiên khách và giỏ món; vòng đời đơn hàng/hóa đơn; KDS cho bếp/bar; ghi nhận thanh toán; nhân sự và phân quyền; dashboard/reporting mức MVP cho chủ/quản lý và quản trị nền tảng. Các tác nhân chính gồm Super Admin, chủ nhà hàng, quản lý, nhân viên phục vụ/POS, bếp/bar và khách hàng quét QR.

Về phạm vi kỹ thuật, khóa luận tập trung vào kiến trúc microservices trong Nx monorepo, với các dịch vụ và client được tổ chức theo ranh giới nghiệp vụ. Các chủ đề kỹ thuật chính gồm service boundary, data ownership, tenant isolation, RBAC, entitlement theo gói, event-driven có chọn lọc, Kafka cho domain event, Redis cho cache/bản chiếu vận hành, WebSocket cho cập nhật gần thời gian thực, Keycloak/JWT/OIDC cho nhân viên và quản trị, token phiên cho khách QR, SePay/VietQR cho thanh toán và idempotency cho gửi đơn/webhook.

Về phạm vi đánh giá, khóa luận sử dụng ma trận truy vết yêu cầu, kiểm thử đơn vị, kiểm thử hợp đồng, kiểm thử tích hợp, phân tích kiến trúc, sơ đồ, ảnh minh họa client, kết quả kiểm thử và hiện vật triển khai production. Việc triển khai lên môi trường production là một nhánh công việc chắc chắn của quá trình hoàn thiện khóa luận; các hiện vật như URL/domain, HTTPS, kiểm tra sức khỏe, kiểm thử khói, đường gọi lại webhook và nhật ký rút gọn được dùng để củng cố phần minh chứng khi đã thu thập. Điều này giúp phần đánh giá nối được từ môi trường phát triển nội bộ đến minh chứng vận hành trên môi trường production, nhưng không biến các hiện vật minh họa thành thay thế cho kiểm thử chức năng hoặc kiểm chứng kiến trúc.

Một số nội dung sản phẩm và vận hành nâng cao không thuộc kết quả chính của khóa luận. Nhóm này gồm đo kiểm tải lớn, p95/p99 latency, throughput, kiểm thử tải/áp lực so sánh với monolith, high availability cấp vận hành thật, chaos engineering, native mobile app cho iOS/Android, BI/AI analytics nâng cao, OLAP/data warehouse, CRM/loyalty, HRM/payroll/shift scheduling/time attendance, inventory/BOM nâng cao, thuế/kế toán, tích hợp kế toán bên thứ ba, full offline write queue, background sync và conflict resolver cho mọi thao tác, email/SMS/push notification đầy đủ, nhiều ngân hàng thanh toán đồng thời, refund/partial refund phức tạp và tích hợp sâu với nền tảng giao đồ ăn bên thứ ba. Các nội dung này là hướng mở rộng tự nhiên của sản phẩm, nhưng cần thiết kế, triển khai và bằng chứng riêng nếu đưa vào phạm vi đánh giá chính.

Ranh giới phạm vi này giúp các chương sau tập trung vào đóng góp cốt lõi: một nền tảng SaaS POS có QR ordering, KDS, payment, quản trị tenant, phân quyền và service boundary rõ. Những giới hạn về tải lớn, sản phẩm nâng cao hoặc vận hành cấp doanh nghiệp không làm giảm mục tiêu của đề tài; chúng xác định phạm vi đánh giá để khóa luận không đưa ra kết luận vượt quá bằng chứng.

\section{Phương pháp thực hiện}

Phương pháp thực hiện khóa luận gồm năm bước chính.

Bước thứ nhất là khảo sát bối cảnh và tài liệu nền. Khóa luận sử dụng nguồn thị trường/sản phẩm để mô tả bối cảnh F\&B, POS, QR ordering và thanh toán QR/VietQR; đồng thời sử dụng chuẩn, RFC, tài liệu chính thức, sách và bài viết kỹ thuật uy tín để trình bày cơ sở lý thuyết về SaaS, multi-tenancy, microservices, Kafka, WebSocket, JWT/OIDC, idempotency và Saga. Các nguồn thị trường được dùng đúng vai trò bối cảnh/related systems, không thay thế định nghĩa học thuật hoặc bằng chứng triển khai của QRTable.

Bước thứ hai là phân tích yêu cầu. Từ bối cảnh vận hành F\&B, khóa luận xác định các tác nhân, use case, yêu cầu chức năng, yêu cầu phi chức năng và máy trạng thái nghiệp vụ. Phân tích này làm rõ mỗi nhóm người dùng cần thao tác gì, dữ liệu nào được tạo hoặc thay đổi, ràng buộc tenant/quyền/gói nào phải được kiểm tra và luồng nghiệp vụ nào cần được đánh giá ở các chương sau.

Bước thứ ba là thiết kế kiến trúc. Hệ thống được tổ chức theo các miền nghiệp vụ như Catalog, Order, Kitchen, Payment, SaaS, User-Access và Authorizer. Thiết kế tập trung vào service boundary, data ownership, mô hình giao tiếp đồng bộ/bất đồng bộ, Kafka topic contract, Redis ownership, WebSocket hint/refetch, bảo mật/RBAC/tenant isolation và các quyết định công nghệ. Các sơ đồ kiến trúc, bảng ownership và bảng trade-off được dùng để giải thích vì sao mỗi quyết định phù hợp với bài toán QRTable.

Bước thứ tư là triển khai hệ thống. QRTable được hiện thực trong Nx monorepo với backend, frontend và các dịch vụ phụ trợ. Backend dùng các dịch vụ tách theo miền nghiệp vụ; client gồm Customer PWA cho khách quét QR và Management App cho chủ/quản lý/nhân viên/quản trị nền tảng. Các luồng cốt lõi được trình bày theo hướng evidence-driven: mục tiêu nghiệp vụ, bất biến cần giữ, dịch vụ sở hữu dữ liệu, mã nguồn/kiểm thử liên quan, sơ đồ trình tự và ảnh minh họa client.

Bước thứ năm là kiểm chứng và đánh giá. Khóa luận kết hợp nhiều lớp bằng chứng: kiểm thử tự động, kiểm thử có điều kiện, kiểm chứng kiến trúc, ma trận truy vết, demo/screenshot và hiện vật triển khai production. Việc đánh giá được tổ chức theo loại bằng chứng để phân biệt điều đã kiểm chứng, điều được hỗ trợ bởi thiết kế và điều cần phát triển thêm. Cách tiếp cận này phù hợp với một khóa luận xây dựng hệ thống phần mềm, vì nó vừa chứng minh được kết quả triển khai, vừa giữ ranh giới học thuật cho các kết luận về hiệu năng, khả năng mở rộng và vận hành.

\section{Đóng góp của đề tài}

Đóng góp thứ nhất của đề tài là mô hình hóa bài toán SaaS POS tích hợp QR ordering cho bối cảnh F\&B. Khóa luận không chỉ xem đặt món qua QR như một trang thực đơn điện tử, mà đặt nó trong chuỗi vận hành gồm bàn, phiên, giỏ, đơn, bếp, hóa đơn, thanh toán, báo cáo và quản trị tenant. Cách mô hình hóa này giúp làm rõ các bất biến nghiệp vụ cần được bảo vệ khi nhiều tác nhân và nhiều thiết bị cùng tham gia một phiên phục vụ.

Đóng góp thứ hai là thiết kế kiến trúc microservices với ranh giới nghiệp vụ và quyền sở hữu dữ liệu rõ. Các miền Catalog, Order, Kitchen, Payment, SaaS, User-Access và Authorizer được phân tách theo trách nhiệm, không chia theo tầng kỹ thuật thuần túy. Thiết kế này tạo cơ sở để giải thích các quyết định về database-per-service, giao tiếp liên dịch vụ, Kafka event, Redis cache/bản chiếu vận hành, WebSocket cập nhật gần thời gian thực và guard chain ở BFF.

Đóng góp thứ ba là hiện thực các luồng vận hành cốt lõi của QRTable. Các luồng gồm khách quét QR và đặt món, nhân viên xác nhận đơn và bảo toàn tồn kho, KDS nhận phiếu bếp, thanh toán cash/VietQR/SePay, khởi tạo tenant theo SaaS Onboarding Mini-Saga, dashboard/reporting theo gói và phân tích nền tảng cho Super Admin. Những luồng này cho thấy hệ thống không chỉ dừng ở thiết kế, mà có mã nguồn, kiểm thử, sơ đồ và minh chứng giao diện tương ứng.

Đóng góp thứ tư là áp dụng các cơ chế quan trọng của hệ thống phân tán vào bối cảnh POS F\&B nhiều tenant. QRTable sử dụng tenant isolation để bảo vệ dữ liệu nhà hàng, RBAC và entitlement để kiểm soát quyền thao tác, idempotency để chống gửi đơn/webhook trùng, Saga đại diện để phối hợp tác dụng phụ liên dịch vụ và WebSocket theo hướng hint/refetch để tránh biến kênh thời gian thực thành nguồn sự thật nghiệp vụ.

Đóng góp thứ năm là xây dựng bộ tài liệu và quy trình đánh giá có khả năng truy vết. Khóa luận liên kết yêu cầu, thiết kế, triển khai, kiểm thử, sơ đồ, screenshot và hiện vật triển khai production thành một mạch bằng chứng. Điều này giúp người đọc không chỉ thấy QRTable có chức năng gì, mà còn hiểu chức năng đó được thiết kế theo ranh giới nào, được kiểm chứng bằng loại bằng chứng nào và cần mở rộng theo hướng nào.

\section{Cấu trúc khóa luận}

Khóa luận được tổ chức thành bảy chương chính.

Chương 1 trình bày bối cảnh, lý do chọn đề tài, phát biểu bài toán, mục tiêu, đối tượng, phạm vi, phương pháp thực hiện và đóng góp của QRTable.

Chương 2 trình bày cơ sở lý thuyết và các công trình/hệ thống liên quan. Nội dung gồm POS trong F\&B, đặt món qua mã QR, SaaS và multi-tenancy, microservices, event-driven architecture, consistency/idempotency, WebSocket/realtime và security trong SaaS POS.

Chương 3 chuyển bối cảnh vận hành F\&B thành yêu cầu hệ thống cho QRTable. Chương này xác định tác nhân, use case, yêu cầu chức năng, yêu cầu phi chức năng, máy trạng thái chính và phạm vi yêu cầu.

Chương 4 trình bày thiết kế kiến trúc và quyết định công nghệ. Chương này tập trung vào service boundary, data ownership, kiến trúc Nx monorepo, giao tiếp liên dịch vụ, Kafka, Redis, bảo mật/RBAC/tenant isolation và tích hợp thanh toán.

Chương 5 trình bày quá trình hiện thực các luồng vận hành cốt lõi. Các luồng chính gồm QR session/cart/order, xác nhận đơn và tồn kho, KDS, thanh toán, SaaS onboarding, dashboard/reporting và đóng gói triển khai production.

Chương 6 kiểm chứng và đánh giá hệ thống. Chương này sử dụng ma trận truy vết, kiểm thử, phân tích kiến trúc, minh chứng demo và hiện vật triển khai production để phân loại bằng chứng cho từng nhóm kết luận.

Chương 7 tổng kết kết quả đạt được, đóng góp của đề tài, các hạn chế còn lại và hướng phát triển. Phần này nối lại mục tiêu ban đầu với kết quả triển khai/đánh giá và xác định các hướng mở rộng phù hợp cho QRTable sau phạm vi khóa luận.
